\chapter{Conclusion}

Computer architecture requires learners to connect several representations of instruction execution that are often introduced separately: instruction formats, registers, control decisions, datapath structure, memory behaviour, and changes to architectural state. As established in Chapter~2, difficulty can persist even when the individual components are familiar, because the learner must still turn those partial representations into a coherent account of what the processor is doing \cite{porter2013evaluating,yehezkel2007contribution}. This dissertation addressed that narrower educational problem through a virtual reality learning environment focused on selected single-cycle MIPS datapath concepts and positioned throughout as a supplement to conventional teaching.

The resulting prototype was designed around guided instruction tracing rather than around a literal three-dimensional reconstruction of the processor. Selected datapath relationships were redistributed across six learner-facing phases---Fetch, Decode, Execute, Memory, Write-Back, and Program Counter Update---and supported through physical instruction handling, register interaction, visible value transfer, local validation, and progressively reduced guidance across Learning, Practice, and Test modes. The prototype was then evaluated through the exploratory participant study described in Chapter~3, with the findings reported in Chapter~6 and interpreted in Chapter~7.

\section{Main Contribution and Research Question}

The main contribution lies in the way the single-cycle datapath was translated for teaching. Hardware behaviour that is compact, partly hidden, and effectively concurrent was reorganised into a sequential reasoning path that learners could see, follow, act upon, and revisit. The aim was not to reproduce processor timing literally, but to preserve the relationships that a learner must understand while making those relationships more legible during instruction. The participant evidence provides some support for this design choice: visualisation, mental-model formation, direct interaction, and local guidance were repeatedly identified as useful, while the study also exposed where the translation remained weaker, particularly around some control-flow, memory, and phase-boundary concepts.

The research question asked:

\begin{quote}
\textit{To what extent can an immersive virtual reality learning environment support students' understanding of selected single-cycle MIPS datapath concepts as a supplementary educational tool?}
\end{quote}

Within the limits of the evaluation, the answer is qualified but positive. The prototype appears capable of supporting understanding of selected single-cycle MIPS datapath concepts when used as a supplementary learning environment, with its strongest value lying in the visual and interactive organisation of instruction execution into a learner-facing sequence. Participants consistently reported that the visualisations helped them form a mental model and made abstract processes easier to understand, while post-session knowledge performance was generally strong across several targeted concepts. All participants also progressed from Learning into Practice, providing limited evidence that the environment could move beyond fully guided explanation into an initial reduced-support condition.

The strength of that conclusion remains deliberately limited. The evaluation involved a small sample, no non-VR comparison condition, and no delayed retention measure. The tailored refresher delivered between the baseline probe and the headset session also means that descriptive baseline/post-session differences reflect the overall learning sequence rather than an isolated effect of VR. Progression into Test mode was limited, some concepts remained fragile, and physical discomfort or interface friction affected parts of the participant experience. The findings are therefore suggestive rather than definitive. They support further development of this particular design approach, but they do not demonstrate that VR is generally superior to lectures, diagrams, tutorials, simulators, or other established forms of computer-architecture teaching.

This distinction also defines the educational role of the prototype more clearly. Its strongest use is not as a replacement for earlier instruction, but as an additional representation after learners have already encountered the terminology and structure of the datapath. In that position, the environment can slow instruction execution down, make selected relationships visible, and provide another route from recognising the components of the processor to explaining how they participate in a complete instruction walkthrough. This supplementary role is consistent with the wider cautious positioning of VR in computer science education, where immersive environments are most defensible when they clarify difficult conceptual or process-oriented material without assuming that immersion itself guarantees stronger learning \cite{janani2026virtual,dengel2020important}.

\section{Future Work}

The most immediate future work follows from the participant study itself. A stronger evaluation would use a larger and more clearly defined participant sample, ideally involving students currently taking a computer architecture module, together with a comparative condition using an alternative teaching format. A more standardised lesson sequence would also make progression across Learning, Practice, and Test easier to interpret, particularly if students encountered a common set of instructions before later variation was introduced. Delayed assessment would then be needed to determine whether the mental-model and concept support reported immediately after the session persisted over time. The observational component could likewise be strengthened through a predefined protocol that distinguishes conceptual assistance from navigation or interface assistance.

The second priority is refinement of the current VR experience. Comfort should be addressed directly through greater control over locomotion speed, reduced reliance on sustained smooth movement, clearer alternatives such as teleportation, and route design that limits unnecessary travel between stations. These issues matter because the environment cannot function well as a teaching aid if physical strain competes with the material itself. The interface should also be simplified where the study identified text density, scrolling, uncertain task order, or unclear interaction targets. Particular attention should be given to preserving stable interaction conventions as guidance is reduced, so that Practice and Test modes assess architectural reasoning rather than memory for the interface. Concepts that remained weaker in the study, including branch-to-PC reasoning, memory-transfer direction, and the relationship between neighbouring phases such as Fetch and Decode, should also receive stronger instructional reinforcement.

Beyond those study-led refinements, broader technical expansion becomes a longer-term priority. The supported instruction set could first be widened so that more arithmetic, memory, and control-flow cases can be explored without changing the underlying lesson model. A later version could then investigate whether the same pedagogical-serialization approach remains useful for multicycle execution or pipelined behaviour, following the progression students would normally encounter as a computer architecture course moves from basic instruction execution toward stalls, hazards, and overlapping instruction flow. These extensions would not simply add more content; they would test whether the central design principle developed here continues to hold once the processor can no longer be explained through one bounded single-cycle walkthrough.

The dissertation therefore closes with a narrower claim than one that VR broadly improves computer-architecture education, but a more defensible one. A carefully scoped immersive environment can provide meaningful supplementary support when it is used to make hidden computational relationships visible, ordered, and interactive. The prototype and study provide an initial demonstration of that approach for selected single-cycle MIPS datapath concepts, while also showing that its educational value depends on clarity, comfort, appropriately reduced guidance, and careful control of scope. Further comparative and longitudinal work is needed before broader conclusions can be drawn, but the design direction itself has enough support to justify that next stage.